\documentclass[11pt]{article}
\usepackage{amsmath, amssymb}
\usepackage{geometry}
\usepackage{array}
\usepackage{booktabs}
\geometry{margin=0.7in}
\setlength{\parindent}{0pt}
\setlength{\parskip}{0.35em}

\begin{document}

\begin{center}
{\LARGE Trigonometry Cheatsheet}\\[0.4em]
{\large Fundamental Review}
\end{center}

\section*{Angle Basics}
\[
180^\circ=\pi \text{ radians},\qquad
360^\circ=2\pi \text{ radians},\qquad
\theta^\circ=\theta\cdot \frac{\pi}{180},\qquad
\theta \text{ rad}=\theta\cdot \frac{180}{\pi}^\circ
\]

\begin{center}
\begin{tabular}{ccccccccc}
\toprule
Degrees & $0^\circ$ & $30^\circ$ & $45^\circ$ & $60^\circ$ & $90^\circ$ & $180^\circ$ & $270^\circ$ & $360^\circ$ \\
Radians & $0$ & $\frac{\pi}{6}$ & $\frac{\pi}{4}$ & $\frac{\pi}{3}$ & $\frac{\pi}{2}$ & $\pi$ & $\frac{3\pi}{2}$ & $2\pi$ \\
\bottomrule
\end{tabular}
\end{center}

\section*{Unit Circle and Reference Angles}
Special reference angles:
\[
\frac{\pi}{6},\qquad \frac{\pi}{4},\qquad \frac{\pi}{3}
\]

For those angles:
\[
\sin\theta:\ \frac12,\ \frac{\sqrt2}{2},\ \frac{\sqrt3}{2}
\qquad
\cos\theta:\ \frac{\sqrt3}{2},\ \frac{\sqrt2}{2},\ \frac12
\qquad
\tan\theta:\ \frac{\sqrt3}{3},\ 1,\ \sqrt3
\]

Quadrant positive sign rule:
\[
\text{QI: all},\qquad \text{QII: }\sin,\qquad \text{QIII: }\tan,\qquad \text{QIV: }\cos
\]

Reference angle formulas:
\[
\theta_{\text{ref}}=
\begin{cases}
\theta & \text{in QI}\\
\pi-\theta & \text{in QII}\\
\theta-\pi \text{ or } \theta+\pi & \text{in QIII}\\
-\theta \text{ or } 2\pi-\theta & \text{in QIV}
\end{cases}
\]

\section*{Right Triangle Trig}
SOH-CAH-TOA:
\[
\sin\theta=\frac{\text{opp}}{\text{hyp}},\qquad
\cos\theta=\frac{\text{adj}}{\text{hyp}},\qquad
\tan\theta=\frac{\text{opp}}{\text{adj}}=\frac{\sin\theta}{\cos\theta}
\]

Reciprocal functions:
\[
\csc\theta=\frac{1}{\sin\theta},\qquad
\sec\theta=\frac{1}{\cos\theta},\qquad
\cot\theta=\frac{1}{\tan\theta}=\frac{\cos\theta}{\sin\theta}
\]

\section*{Core Identities}
Pythagorean identities:
\[
\sin^2\theta+\cos^2\theta=1
\]
\[
1+\tan^2\theta=\sec^2\theta
\]
\[
1+\cot^2\theta=\csc^2\theta
\]

\section*{Cofunction Identities}
\[
\sin\left(\frac{\pi}{2}-\theta\right)=\cos\theta,\qquad
\cos\left(\frac{\pi}{2}-\theta\right)=\sin\theta
\]
\[
\tan\left(\frac{\pi}{2}-\theta\right)=\cot\theta,\qquad
\cot\left(\frac{\pi}{2}-\theta\right)=\tan\theta
\]
\[
\sec\left(\frac{\pi}{2}-\theta\right)=\csc\theta,\qquad
\csc\left(\frac{\pi}{2}-\theta\right)=\sec\theta
\]

\section*{Even-Oddness and Symmetry}
\begin{center}
\begin{tabular}{cccc}
\toprule
Function & Parity & Algebra & Symmetry \\
\midrule
$\sin x$ & odd & $\sin(-x)=-\sin(x)$ & origin \\
$\cos x$ & even & $\cos(-x)=\cos(x)$ & $y$-axis \\
$\tan x$ & odd & $\tan(-x)=-\tan(x)$ & origin \\
$\csc x$ & odd & $\csc(-x)=-\csc(x)$ & origin \\
$\sec x$ & even & $\sec(-x)=\sec(x)$ & $y$-axis \\
$\cot x$ & odd & $\cot(-x)=-\cot(x)$ & origin \\
\bottomrule
\end{tabular}
\end{center}

\section*{Sum and Difference Formulas}
\[
\sin(\alpha+\beta)=\sin\alpha\cos\beta+\cos\alpha\sin\beta
\]
\[
\sin(\alpha-\beta)=\sin\alpha\cos\beta-\cos\alpha\sin\beta
\]
\[
\cos(\alpha+\beta)=\cos\alpha\cos\beta-\sin\alpha\sin\beta
\]
\[
\cos(\alpha-\beta)=\cos\alpha\cos\beta+\sin\alpha\sin\beta
\]
\[
\tan(\alpha+\beta)=\frac{\tan\alpha+\tan\beta}{1-\tan\alpha\tan\beta}
\]
\[
\tan(\alpha-\beta)=\frac{\tan\alpha-\tan\beta}{1+\tan\alpha\tan\beta}
\]

\section*{Double-Angle Formulas}
\[
\sin(2\theta)=2\sin\theta\cos\theta
\]
\[
\cos(2\theta)=
\begin{cases}
\cos^2\theta-\sin^2\theta\\
2\cos^2\theta-1\\
1-2\sin^2\theta
\end{cases}
\]
\[
\tan(2\theta)=\frac{2\tan\theta}{1-\tan^2\theta}
\]

\section*{Half-Angle Formulas}
\[
\sin\left(\frac{\theta}{2}\right)=\pm\sqrt{\frac{1-\cos\theta}{2}}
\]
\[
\cos\left(\frac{\theta}{2}\right)=\pm\sqrt{\frac{1+\cos\theta}{2}}
\]
\[
\tan\left(\frac{\theta}{2}\right)=
\begin{cases}
\dfrac{\sin\theta}{1+\cos\theta}\\[0.6em]
\dfrac{1-\cos\theta}{\sin\theta}\\[0.6em]
\csc\theta-\cot\theta
\end{cases}
\]
The sign depends on the quadrant of $\frac{\theta}{2}$.

\section*{Product-to-Sum Formulas}
\[
\sin\alpha\sin\beta=\frac12\bigl[\cos(\alpha-\beta)-\cos(\alpha+\beta)\bigr]
\]
\[
\cos\alpha\cos\beta=\frac12\bigl[\cos(\alpha-\beta)+\cos(\alpha+\beta)\bigr]
\]
\[
\sin\alpha\cos\beta=\frac12\bigl[\sin(\alpha+\beta)+\sin(\alpha-\beta)\bigr]
\]
\[
\cos\alpha\sin\beta=\frac12\bigl[\sin(\alpha+\beta)-\sin(\alpha-\beta)\bigr]
\]

\section*{Triangles}
Standard labeling: side $a$ opposite angle $A$, side $b$ opposite angle $B$, side $c$ opposite angle $C$.

Area with two sides and included angle:
\[
K=\frac12 ab\sin C=\frac12 bc\sin A=\frac12 ca\sin B
\]

Law of Sines:
\[
\frac{a}{\sin A}=\frac{b}{\sin B}=\frac{c}{\sin C}
\]

SSA ambiguous case with known $A,a,b$:
\[
h=b\sin A
\]
\begin{itemize}
  \item If $a<h$, there are $0$ triangles.
  \item If $a=h$, there is $1$ right triangle.
  \item If $h<a<b$, there are $2$ triangles.
  \item If $a\ge b$, there is $1$ triangle.
\end{itemize}

Law of Cosines:
\[
c^2=a^2+b^2-2ab\cos C
\]
\[
a^2=b^2+c^2-2bc\cos A,\qquad
b^2=a^2+c^2-2ac\cos B
\]

Use Law of Cosines for:
\begin{itemize}
  \item SAS: two sides and included angle
  \item SSS: all three sides known
\end{itemize}

Heron's Formula:
\[
s=\frac{a+b+c}{2},\qquad
K=\sqrt{s(s-a)(s-b)(s-c)}
\]

Brahmagupta's Formula for a cyclic quadrilateral (inscribed in a circle) with side lengths $a,b,c,d$:
\[
s=\frac{a+b+c+d}{2},\qquad
K=\sqrt{(s-a)(s-b)(s-c)(s-d)}
\]

\section*{Bearings}
True bearings are measured clockwise from north.

\[
\text{N }35^\circ\text{ E} \Rightarrow 35^\circ \text{ east of north}
\]
\[
\text{S }20^\circ\text{ W} \Rightarrow 20^\circ \text{ west of south}
\]

Common conversion idea:
\begin{itemize}
  \item Bearing $060^\circ$ means $60^\circ$ clockwise from north.
  \item Bearing $210^\circ$ means $210^\circ$ clockwise from north.
  \item Draw a north-east-south-west diagram first, then mark the angle carefully.
\end{itemize}

For triangle word problems with bearings:
\begin{itemize}
  \item convert each direction to a standard diagram angle
  \item find the included angle between paths
  \item use Law of Sines or Law of Cosines as needed
\end{itemize}

\end{document}
